Całość systemu składa się na w pełni funkcjonalną aplikację internetową, będącą serwisem z ogłoszeniami o sprzedaży samochodów.
Użytkownicy mają możliwość wystawiania swoich pojazdów na sprzedaż oraz przeglądania istniejących ofert.
Aplikacja opatrzona jest możliwością publikacji komentarzy oraz wystawiania ocen pod ogłoszeniami.
W celu zapewnienia bezpieczeństwa społeczności, do rozwiązania włączony jest system moderowania treściami dostępnymi w serwisie.

\section{Autentykacja i autoryzacja użytkowników}

Proces autentykacji i autoryzacji użytkowników zapewnia w sposób kontrolowany dostęp do poszczególnych funkcjonalności serwisu, zgodnie z zasadą najmniejszych uprawnień \autocite{crowdstrike_polp}.
Rejestracja konta w serwisie odbywa się poprzez formularz, gdzie użytkownik podaje podstawowe informacje, takie jak adres e-mail, na który jest wysyłany link aktywacyjny w celu sfinalizowania procesu.
Anonimowi odwiedzający stronę, którzy nie zostali poddani autentykacji, są ograniczeni jedynie do przeglądania ogłoszeń o sprzedaży stworzonych przez innych użytkowników oraz zamieszczonych przez nich opinii.
Po pomyślnym zalogowaniu, przy pomocy danych uwierzytelniających podanych w procesie rejestracji, użytkownicy uzyskują dostęp do rozszerzonej funkcjonalności systemu.
Funkcjonalności te obejmują wystawianie ofert o sprzedaży, dodawanie komentarzy i ocen pod ogłoszeniami oraz zgłaszanie treści naruszających politykę serwisu.
System autoryzacji jest oparty o model ról i~uprawnień, gdzie niektórzy z użytkowników mają możliwość przeglądania nowych zgłoszeń i ewentualnego usuwania objętych nimi treści.

\section{Zamieszczanie ogłoszeń o sprzedaży samochodów}

Zalogowani użytkownicy posiadają możliwość publikacji ofert sprzedaży swoich pojazdów poprzez utworzenie ogłoszenia w dedykowanym panelu w aplikacji.
Kreator ogłoszeń jest zestawem kolejno następujących po sobie formularzy, intuicyjnie kierując użytkownika przez cały proces.
W pierwszym etapie użytkownik podaje parametry pojazdu, gdzie większość informacji podaje się w predefiniowanych polach wyboru, co znacząco usprawnia proces wypełniania formularza.
Kolejnym etapem jest przekazanie metadanych ogłoszenia, takich jak tytuł, opis oferty lub cenę pojazdu.
Ostatni z formularzy umożliwia dodanie zdjęć przedstawiających wygląd pojazdu wewnątrz oraz na zewnątrz.
Jeżeli wszystkie dane zostaną uzupełnione poprawnie, użytkownik może opublikować ogłoszenie.
Na tym etapie oferta będzie poddana weryfikacji przez osoby moderujące stronę i ostatecznie przez nich zatwierdzona pod względem zgodności z polityką strony oraz stanem faktycznym pojazdu.

\section{Dodawanie i ocenianie opinii pod ofertami}

Każde ogłoszenie zawiera dedykowaną sekcję komentarzy, w której użytkownicy dzielą się opiniami dotyczącymi pojazdu, bazując na treści samej oferty, jak również na podstawie własnych doświadczeń.
Poszczególne komentarze są również poddawane ocenie poprzez zastosowanie binarnego systemu ocen opartego na pozytywnych i negatywnych reakcji użytkowników.
Oceny te determinują wagę poszczególnych opinii w sumarycznej ocenie prezentowanej oferty oraz ograniczają wpływ niekonstruktywnych lub nawet szkodliwych komentarzy.
System ten wspiera mniej doświadczonych kupującym na bardziej przemyślane decyzje, dzięki opiniom użytkowników o większym doświadczeniu, a jednocześnie pozwala na sprawniejszą identyfikację nieuczciwych ofert.

\section{Wyszukiwarka ofert w oparciu o reputację sprzedawcy}

Jednym z elementów systemu, dostępnym dla wszystkich użytkowników, jest wyszukiwarka ofert.
Funkcjonalność ta jest dostępna bezpośrednio z poziomu ekranu głównego aplikacji i prezentuje listę skróconych wersji ofert w postaci interaktywnych kafelków.
Głównym filtrem, o który oparte jest wyszukiwanie, jest pole tekstowe, które generuje listę ofert względem dopasowania ich treści do podanej sekwencji słów.
Ponadto, dostępny jest wybór wielu predefiniowanych filtrów pozwalających na dokładne dostosowanie wyszukiwania, na przykład względem marki pojazdu lub przedziału cenowego.
Kolejność wyników bazuje na aktualności ofert, ich ocenie oraz historycznej reputacji ich właściciela.
Wybór konkretnego kafelka skutkuje przekierowaniem użytkownika do osobnej strony, zawierającej szczegółowe informacje dotyczące danej oferty.

\section{System zgłoszeń dla ofert i komentarzy}

Zarówno oferty, jak i komentarze mogą być zgłoszone przez użytkowników aplikacji za pośrednictwem dedykowanego przycisku.
Jego naciśnięcie skutkuje wyświetleniem formularza zgłoszeniowego, gdzie istnieje możliwość wskazania przyczyny zgłoszenia wybranej treści.
System raportowania oparty jest na wyselekcjonowanej grupie użytkowników, którzy pełnią rolę moderatorów treści i~posiadają dostęp do gromadzonych przez serwis zgłoszeń.
Poprzez stworzony dla nich panel administracyjny, mogą dokonywać oni zbiorczego przeglądu zgłoszeń i~podejmować decyzje dotyczące dostępności oznaczonych treści.
Zgłoszenia są prezentowane w postaci listy, gdzie zawarte są informacje o rodzaju objętej nimi treści, użytkowniku dokonującym zgłoszenia oraz dodatkowych informacjach przekazanych w formularzu.

\section{Przechowywanie nieaktualnych ogłoszeń w archiwum}

Kiedy ogłoszenie wygasa, na przykład na skutek zakończonej sprzedaży bądź anulowania oferty, nie jest kompletnie usuwane z systemu, lecz zostaje przeniesione do archiwum.
Ogłoszenia oznaczone jako zarchiwizowane nie są dostępne z poziomu standardowej wyszukiwarki, natomiast pozostają one trwale przypisane do profili użytkowników, które je wystawiły.
Użytkownik zainteresowany daną ofertą ma możliwość przejścia do widoku profilu jej właściciela i zapoznania się z pełną historią ogłoszeń w postaci listy, posortowanej według daty opublikowania od najnowszych.
Zaimplementowane rozwiązanie znacznie zwiększa transparentność sprzedawców, uniemożliwiając manipulowanie historią ofert poprzez usuwanie niepożądanych przez nich wpisów.

\section{Responsywność}

Aplikacja automatycznie dostosowuje się do większości urządzeń konsumenckich, które obsługują popularne przeglądarki.
Wszystkie komponenty aplikacji dynamicznie zmieniają swoje położenie, geometrię i wygląd w zależności od wymiarów ekranu oraz orientacji urządzenia.
Responsywność zapewnia, że wszystkie składowe serwisu są czytelne i wygodne w użytkowaniu niezależnie od platformy, na którym jest uruchamiany.

\section{Dostępność}

Każda z funkcjonalności serwisu zostały zaprojektowane z możliwością obsługi klawiaturowej bez konieczności używania kursora.
Kolory tekstu oraz tła zostały dobrane tak aby spełniać minimalne wymagania kontrastu WCAG \autocite{wcag_guidelines}.
Wszystkie komponenty dostosowują swoją kolorystykę zależnie od wybranego przez użytkownika motywu, jasnego lub ciemnego.
Wszystkie pola formularzy zapewniają podstawowe wsparcie dla czytników ekranowych, a obrazy posiadają teksty alternatywne, które są wyświetlane na wypadek braku możliwości wyświetlenia.
Niezależnie od występujących ograniczeń większość użytkowników posiada możliwości korzystania z serwisu w podobny sposób.

\section{Obsługa błędów}

Podczas korzystania z aplikacji mogą występować błędy mające różne źródła.
Mogą to być niepoprawne formaty danych w przesyłanych formularzach, komunikaty błędów zwracane z serwera lub nawet wewnętrzne błędy samej aplikacji.
Niezależnie od pochodzenia każdy z błędów jest wyłapywany i przetwarzany na komunikat, który jest czytelny dla użytkownika oraz propozycja rozwiązania problemu.
Kontrola nad błędami zapewnia poprawia doświadczenie użytkownika w korzystaniu z serwisu oraz zwiększa wydajność całego systemu.
